\chapter{Równania Riccatiego i sterowanie liniowo-kwadratowe}
\label{ch:lin24-riccati}
\index{równanie Riccatiego@Równanie Riccatiego}
\index{sterowanie liniowo-kwadratowe@Sterowanie liniowo-kwadratowe}
\index{równanie Riccatiego!algebraiczne}
\index{sprzężenie zwrotne@Sprzężenie zwrotne}

\noindent\textbf{Cel rozdziału.}
Twierdzenia Lapunowa z
Rozdziałów~\ref{ch:lin22-sylvester}
i~\ref{ch:lin23-discrete-lyapunov}
rozpoznają stabilność ustalonej macierzy.
Równania Riccatiego pozwalają
równocześnie \emph{wybrać} operator
sprzężenia zwrotnego i uzasadnić
jego optymalność.
Wyprowadzimy je z problemów
minimalizacji całkowitego kosztu
kwadratowego, najpierw w czasie
dyskretnym, potem ciągłym.
Następnie udowodnimy istnienie,
dodatnią określoność i jednoznaczność
rozwiązań stabilizujących
przy jawnych, wystarczających
założeniach. Nie będziemy
zakładać ogólnych twierdzeń
o sterowalności ani wykrywalności,
które należą do odrębnej teorii.

W całym rozdziale \(n,m\ge1\).
Macierze \(A\in\operatorname{Mat}_{n\times n}(\mathbb C)\),
\(B\in\operatorname{Mat}_{n\times m}(\mathbb C)\),
\(Q\in\operatorname{Mat}_{n\times n}(\mathbb C)\)
oraz \(R\in\operatorname{Mat}_{m\times m}(\mathbb C)\)
są ustalone, przy czym
\(Q=Q^*>0\) i \(R=R^*>0\).
Przez \(x^*\) oznaczamy
sprzężenie hermitowskie.
Dla macierzy hermitowskich
\(P_1\le P_2\) oznacza,
że \(P_2-P_1\) jest
nieujemnie określona.
Wszystkie twierdzenia pozostają
ważne nad \(\mathbb R\)
po zastąpieniu sprzężenia
transpozycją.

\section{Minimalizacja formy kwadratowej}
\label{sec:lin24-square}

\begin{lemat}[Uzupełnianie kwadratu]
\label{lem:lin24-square}
Niech \(M=M^*>0\)
będzie macierzą stopnia \(m\)
nad \(\mathbb C\),
\(c\in\mathbb C^m\)
i \(d\in\mathbb R\).
Dla każdego \(u\in\mathbb C^m\)
zachodzi
\[
 u^*Mu+2\operatorname{Re}(u^*c)+d
 =(u+M^{-1}c)^*M(u+M^{-1}c)
    +d-c^*M^{-1}c.
\]
Wyrażenie z lewej strony
ma jedyne minimum
\(d-c^*M^{-1}c\)
osiągane w \(u=-M^{-1}c\).
\end{lemat}
\begin{proof}
Dodatnia określoność \(M\)
implikuje jej odwracalność
i hermitowskość \(M^{-1}\).
Po wymnożeniu prawej strony
otrzymujemy
\(u^*Mu+u^*c+c^*u
 +c^*M^{-1}c+d-c^*M^{-1}c\).
Ponieważ \(c^*u=\overline{u^*c}\),
środkowe składniki
dają \(2\operatorname{Re}(u^*c)\).
Pierwszy składnik po prawej
w sformułowaniu twierdzenia
jest nieujemny i znika dokładnie,
gdy \(u+M^{-1}c=0\).
To daje wartość i jednoznaczność
minimum.
\end{proof}

\begin{uwaga}[Rola założeń]
\label{rem:lin24-hypotheses}
W dalszych twierdzeniach
dodatnia określoność \(R\)
zapewnia ściśle wypukłą
minimalizację po sterowaniu.
Dodatnia określoność \(Q\)
posłuży do wykazania
stabilności zamkniętej pętli.
Założenia te są silniejsze
niż konieczne w ogólnej teorii
sterowania, ale pozwalają
przedstawić kompletne dowody
bez odwołań do nieomówionych
kryteriów sterowalności.
\end{uwaga}

\section{Czas dyskretny: problem o skończonym horyzoncie}
\label{sec:lin24-discrete-finite}

\begin{definicja}[Dyskretny problem LQR]
\label{def:lin24-discrete-cost}
Dla całkowitego \(N\ge0\)
i \(x_0\in\mathbb C^n\)
rozważmy trajektorię
\[
 x_{k+1}=Ax_k+Bu_k,\qquad 0\le k<N,
\]
przy dowolnych sterowaniach
\(u_0,\ldots,u_{N-1}\in\mathbb C^m\).
Koszt określamy wzorem
\[
 J_N(x_0;u_0,\ldots,u_{N-1})
  =\sum_{k=0}^{N-1}
         (x_k^*Qx_k+u_k^*Ru_k).
\]
Przy \(N=0\) sumę uznajemy
za zerową. \emph{Wartością optymalną}
nazywamy
\(V_N(x_0)=\inf J_N\),
gdzie infimum przebiega
wszystkie dopuszczalne ciągi
sterowań.
\end{definicja}

\begin{twierdzenie}[Dyskretne równanie Riccatiego]
\label{thm:lin24-dare-recursion}
Dla macierzy \(A,B,Q,R\)
spełniających założenia rozdziału
zdefiniujmy ciąg macierzy
\(P_N\in\operatorname{Mat}_{n\times n}(\mathbb C)\)
rekurencją \(P_0=0\) oraz
\[
 P_{N+1}
 =Q+A^*P_NA
  -A^*P_NB(R+B^*P_NB)^{-1}B^*P_NA.
\]
Dla każdego \(N\ge0\)
macierz \(P_N\) jest hermitowska
i nieujemnie określona,
a \(R+B^*P_NB>0\).
Ponadto
\[
 V_N(x)=x^*P_Nx
 \qquad(x\in\mathbb C^n),
\]
zaś infimum w definicji
\(V_N\) jest osiągane.
Na pierwszym kroku problemu
o horyzoncie \(N+1\)
optymalny wybór sterowania
jest jednoznaczny i wynosi
\[
 u_0=-K_Nx_0,\qquad
 K_N=(R+B^*P_NB)^{-1}B^*P_NA.
\]
\end{twierdzenie}
\begin{proof}
Przeprowadzimy indukcję
względem \(N\).
Dla \(N=0\) koszt wynosi zero
i \(V_0(x)=0=x^*P_0x\).
Załóżmy, że teza zachodzi
dla pewnego \(N\ge0\).
W problemie o horyzoncie \(N+1\)
pierwszy koszt wynosi
\(x^*Qx+u^*Ru\),
a kolejny stan to \(Ax+Bu\).
Z założenia indukcyjnego
najmniejszy koszt pozostałych
\(N\) kroków jest równy
\((Ax+Bu)^*P_N(Ax+Bu)\).
Dlatego
\[
 V_{N+1}(x)
 =\min_{u\in\mathbb C^m}
 \bigl[x^*Qx+u^*Ru+
       (Ax+Bu)^*P_N(Ax+Bu)\bigr].
\]
W tym kroku zapisaliśmy
minimum, a nie tylko infimum:
za chwilę skonstruujemy
jego jedyny punkt osiągnięcia.

Po rozwinięciu nawiasu
wyrażenie zależne od \(u\)
ma postać
\[
 u^*M_Nu+
  2\operatorname{Re}(u^*B^*P_NAx)
  +x^*(Q+A^*P_NA)x,
\quad M_N=R+B^*P_NB.
\]
Ponieważ \(P_N\ge0\)
i \(R>0\), mamy \(M_N>0\).
Lemat~\ref{lem:lin24-square}
daje jedyne minimum
przy \(u=-M_N^{-1}B^*P_NAx\),
a jego wartość jest równa
\(x^*P_{N+1}x\)
z macierzą określoną w tezie.
Tym samym \(P_{N+1}\)
jest hermitowska
jako suma i różnica
macierzy hermitowskich.
Ponieważ każdy składnik
pierwotnego kosztu jest
nieujemny, otrzymana wartość
optymalna jest nieujemna:
\(x^*P_{N+1}x\ge0\)
dla każdego \(x\).
Zatem \(P_{N+1}\ge0\).
Łącząc wybrane pierwsze
sterowanie z osiągającymi
minimum sterowaniami pozostałego
horyzontu, otrzymujemy
osiągnięcie infimum.
Indukcja jest zakończona.
\end{proof}

\begin{stwierdzenie}[Monotoniczność wartości optymalnych]
\label{prop:lin24-discrete-mono}
Przy założeniach
Twierdzenia~\ref{thm:lin24-dare-recursion}
dla każdego \(N\ge0\)
zachodzi \(P_N\le P_{N+1}\).
Jeżeli dodatkowo \(A\)
jest macierzą Schura
w sensie
Definicji~\ref{def:lin23-schur},
to istnieje macierz \(W=W^*>0\)
taka, że \(P_N\le W\)
dla wszystkich \(N\ge0\).
\end{stwierdzenie}
\begin{proof}
Dla dowolnego sterowania
w problemie o horyzoncie \(N+1\)
jego pierwsze \(N\) kroków
tworzy dopuszczalne sterowanie
dla horyzontu \(N\).
Wszystkie koszty etapowe
są nieujemne, więc koszt
o dłuższym horyzoncie
nie może być mniejszy
od wartości optymalnej
dla krótszego horyzontu.
Biorąc infimum, otrzymujemy
\(V_{N+1}(x)\ge V_N(x)\)
dla każdego \(x\).
Z wyrażenia
\(V_N(x)=x^*P_Nx\)
wynika \(P_N\le P_{N+1}\).

Załóżmy teraz \(\rho(A)<1\).
Z dyskretnego twierdzenia
Lapunowa
\ref{thm:lin23-lyapunov}
istnieje dodatnio określona
macierz
\[
 W=\sum_{k=0}^{\infty}(A^*)^kQA^k.
\]
Wybierając sterowanie
\(u_k=0\) dla każdego \(k\),
otrzymujemy trajektorię
\(x_k=A^kx\) i koszt
\(\sum_{k=0}^{N-1}x^*(A^*)^kQA^kx
 \le x^*Wx\).
Ponieważ \(V_N(x)\)
jest minimum kosztu,
\(x^*P_Nx\le x^*Wx\)
dla każdego \(x\),
czyli \(P_N\le W\).
\end{proof}

\section{Algebraiczne równanie Riccatiego w czasie dyskretnym}
\label{sec:lin24-dare}

\begin{definicja}[Dyskretne algebraiczne równanie Riccatiego]
\label{def:lin24-dare}
Dla \(A,B,Q,R\) jak na początku
rozdziału \emph{dyskretnym
algebraicznym równaniem Riccatiego}
(DARE) nazywamy
\[
 P=Q+A^*PA-
 A^*PB(R+B^*PB)^{-1}B^*PA,
\]
gdzie szukana jest macierz
hermitowska \(P\),
dla której \(R+B^*PB\)
jest odwracalna.
Gdy \(P\ge0\), macierz ta
jest automatycznie
dodatnio określona.
Dla takiego rozwiązania
definiujemy
\[
 K=(R+B^*PB)^{-1}B^*PA,
 \qquad A_{\mathrm{cl}}=A-BK.
\]
Rozwiązanie nazywamy
\emph{stabilizującym},
jeżeli \(A_{\mathrm{cl}}\)
jest macierzą Schura.
\end{definicja}

\begin{twierdzenie}[Istnienie i jednoznaczność stabilizującego DARE]
\label{thm:lin24-dare-existence}
Niech \(A\in\operatorname{Mat}_{n\times n}(\mathbb C)\)
będzie macierzą Schura,
\(B\in\operatorname{Mat}_{n\times m}(\mathbb C)\)
dowolną macierzą,
a \(Q=Q^*>0\) i \(R=R^*>0\).
Wówczas ciąg \(P_N\)
z Twierdzenia~\ref{thm:lin24-dare-recursion}
ma granicę \(P=P^*>0\).
Macierz \(P\) spełnia DARE,
a sprzężenie zwrotne
\(K=(R+B^*PB)^{-1}B^*PA\)
czyni \(A-BK\) macierzą Schura.

Wśród wszystkich
hermitowskich, nieujemnie
określonych rozwiązań DARE
macierz \(P\) jest jedyna.
Jest ona także jedynym
\emph{hermitowskim rozwiązaniem
stabilizującym}, nawet jeżeli
nie zakłada się z góry jego
nieujemnej określoności.
\end{twierdzenie}
\begin{proof}
Z
Twierdzenia~\ref{thm:lin24-dare-recursion}
i Stwierdzenia~\ref{prop:lin24-discrete-mono}
wynika
\(0=P_0\le P_1\le\cdots\le W\),
gdzie \(W=W^*>0\)
nie zależy od \(N\).
Dla każdego ustalonego
\(z\in\mathbb C^n\)
ciąg \(z^*P_Nz\) jest zatem
niemalejący i ograniczony,
więc ma granicę.
W szczególności zbieżne są
wszystkie elementy diagonalne
\(e_i^*P_Ne_i\).
Wybierając kolejno
\(z=e_i+e_j\) i
\(z=e_i+ie_j\),
odtwarzamy z granic
części rzeczywiste i urojone
elementów pozadiagonalnych
dzięki hermitowskości \(P_N\).
Zatem wszystkie współczynniki
\(P_N\) są zbieżne.
Ich granica \(P\)
jest hermitowska i spełnia
\(0\le P\le W\).
Ponieważ \(P_1=Q>0\)
i \(P_N\) jest niemalejący,
mamy nawet \(P\ge Q>0\).

Dla wszystkich \(N\)
macierz \(M_N=R+B^*P_NB\)
spełnia \(M_N\ge R>0\).
Jej granica
\(M=R+B^*PB>0\)
jest odwracalna.
Odwracanie macierzy
jest ciągłe na zbiorze
macierzy odwracalnych:
można to sprawdzić
z tożsamości
\(M_N^{-1}-M^{-1}
 =M_N^{-1}(M-M_N)M^{-1}\),
przy czym odwrotności \(M_N^{-1}\)
mają jednostajnie ograniczone
normy z nierówności \(M_N\ge R\).
Przechodząc do granicy
w rekurencji dla \(P_N\),
otrzymujemy DARE.

Oznaczmy
\(K=M^{-1}B^*PA\).
Uzupełnienie kwadratu
z Lematu~\ref{lem:lin24-square},
wstawione do równania DARE,
daje tożsamość
\[
 P=Q+K^*RK+(A-BK)^*P(A-BK).
\]
Dla sprawdzenia wystarczy
rozwinąć prawą stronę
i użyć \(MK=B^*PA\)
oraz \(K^*M=A^*PB\).
Stąd
\[
 A_{\mathrm{cl}}^*PA_{\mathrm{cl}}-P
 =-(Q+K^*RK),
 \qquad A_{\mathrm{cl}}=A-BK.
\]
Macierze \(P\) i
\(Q+K^*RK\) są dodatnio
określone, dlatego
dyskretne twierdzenie Lapunowa
\ref{thm:lin23-lyapunov}
implikuje \(\rho(A_{\mathrm{cl}})<1\).

Zbadamy jednoznaczność.
Niech \(\widetilde P=\widetilde P^*\)
będzie innym rozwiązaniem
DARE, dla którego
\(\widetilde M=R+B^*\widetilde PB>0\),
a macierz
\(\widetilde A_{\mathrm{cl}}
 =A-B\widetilde K\),
gdzie
\(\widetilde K=\widetilde M^{-1}B^*\widetilde PA\),
jest macierzą Schura.
Dla dowolnych \(x,u\)
równanie DARE po
uzupełnieniu kwadratu daje
\[
 x^*Qx+u^*Ru+
 (Ax+Bu)^*\widetilde P(Ax+Bu)
 -x^*\widetilde Px
 =(u+\widetilde Kx)^*
        \widetilde M(u+\widetilde Kx).
\]
Tożsamość jest prawdziwa
bez założenia
\(\widetilde P\ge0\):
wystarcza hermitowskość
\(\widetilde P\)
i odwracalność
\(\widetilde M>0\).
Wybierzmy sterowanie
\(u_k=-Kx_k\).
Ponieważ \(A-BK\)
jest macierzą Schura,
trajektoria \(x_k\)
dąży do zera
na mocy Lematu~\ref{lem:lin23-powers}.
Sumując ostatnią tożsamość
dla \(k=0,\ldots,N-1\)
i przechodząc do granicy
\(N\to\infty\),
otrzymujemy
\[
 J_\infty(x_0;-Kx)
 =x_0^*\widetilde Px_0+
  \sum_{k=0}^{\infty}
  ((\widetilde K-K)x_k)^*
  \widetilde M((\widetilde K-K)x_k)
 \ge x_0^*\widetilde Px_0.
\]
Szereg jest zbieżny,
bo \(x_k\) maleje geometrycznie.
Z analogicznej tożsamości
dla \(P\) i sprzężenia \(K\)
wynika, że to samo sterowanie
ma koszt
\(J_\infty(x_0;-Kx)=x_0^*Px_0\),
ponieważ jego dodatkowy
wyraz kwadratowy jest zerowy.
Stąd \(P\ge\widetilde P\).
Zamieniając role obu
stabilizujących rozwiązań
i wybierając sterowanie
\(u_k=-\widetilde Kx_k\),
otrzymujemy
\(\widetilde P\ge P\).
Wobec tego \(\widetilde P=P\).

Jeśli natomiast
\(\widetilde P\ge0\)
jest dowolnym hermitowskim
rozwiązaniem DARE,
to \(\widetilde M>0\).
Ta sama tożsamość
\[
 \widetilde P
 =Q+\widetilde K^*R\widetilde K+
   \widetilde A_{\mathrm{cl}}^*
      \widetilde P\widetilde A_{\mathrm{cl}}
\]
daje \(\widetilde P\ge Q>0\).
Twierdzenie~\ref{thm:lin23-lyapunov}
pokazuje, że
\(\widetilde A_{\mathrm{cl}}\)
jest macierzą Schura.
Taka macierz \(\widetilde P\)
jest więc rozwiązaniem
stabilizującym i z już
udowodnionej jednoznaczności
wynika \(\widetilde P=P\).
\end{proof}

\begin{wniosek}[Optymalność sterowania nieskończonego]
\label{cor:lin24-discrete-opt}
Przy założeniach
Twierdzenia~\ref{thm:lin24-dare-existence}
niech \(P\) i \(K\)
będą opisanymi w nim macierzami.
Dla każdego \(x_0\in\mathbb C^n\)
wśród wszystkich ciągów sterowań
\((u_k)_{k\ge0}\)
o skończonym koszcie
\[
 J_\infty(x_0;u)
 =\sum_{k=0}^\infty
      (x_k^*Qx_k+u_k^*Ru_k)
\]
jedyny ciąg minimalizujący
spełnia \(u_k=-Kx_k\),
\(x_{k+1}=(A-BK)x_k\),
i ma koszt
\(J_\infty(x_0;u)=x_0^*Px_0\).
\end{wniosek}
\begin{proof}
Dla rozwiązania DARE
rachunek z dowodu
Twierdzenia~\ref{thm:lin24-dare-existence}
daje, dla każdego \(x,u\),
\[
 x^*Qx+u^*Ru+
 (Ax+Bu)^*P(Ax+Bu)-x^*Px
 =(u+Kx)^*M(u+Kx),
\]
gdzie \(M=R+B^*PB>0\).
Zsumowanie tej tożsamości
dla \(k=0,\ldots,N-1\)
daje
\[
 \sum_{k=0}^{N-1}
       (x_k^*Qx_k+u_k^*Ru_k)
 =x_0^*Px_0-x_N^*Px_N+
   \sum_{k=0}^{N-1}
       (u_k+Kx_k)^*M(u_k+Kx_k).
\]
Jeżeli \(J_\infty(x_0;u)<\infty\),
to dodatnia określoność \(Q\)
implikuje \(\sum_k\|x_k\|_2^2<\infty\),
a zatem \(\|x_N\|_2\to0\).
Nie potrzebujemy tu osobnego
twierdzenia o stabilnym splocie:
każdy z wyrazów nieujemnego
zbieżnego szeregu \(\|x_k\|_2^2\)
musi dążyć do zera.
Dlatego po przejściu do granicy
otrzymujemy
\[
 J_\infty(x_0;u)
 =x_0^*Px_0+
   \sum_{k=0}^\infty
       (u_k+Kx_k)^*M(u_k+Kx_k)
 \ge x_0^*Px_0.
\]
Nierówność jest ostra, jeśli
choć jedno \(u_k+Kx_k\ne0\),
z dodatniej określoności \(M\).
Dla sprzężenia \(u_k=-Kx_k\)
zamknięta macierz \(A-BK\)
jest stabilna,
więc \(\|x_k\|\) i \(\|u_k\|\)
maleją geometrycznie.
Jego koszt jest skończony
i wszystkie dodatkowe wyrazy
kwadratowe znikają.
Jest ono zatem jedynym
ciągiem minimalizującym.
\end{proof}

\begin{przyklad}[Skalarne DARE, dwa rozwiązania]
\label{prz:lin24-dare-scalar}
Niech \(n=m=1\),
\(A=1/2\), \(B=1\),
\(Q=7/8\) i \(R=1\).
Macierz \(A\) jest stabilna.
Rekurencja z
Twierdzenia~\ref{thm:lin24-dare-recursion}
ma w tym przypadku postać
\[
 p_0=0,\qquad
 p_{N+1}=\frac78+
         \frac{p_N}{4(1+p_N)}.
\]
W szczególności
\(p_1=7/8\) oraz
\(p_2=7/8+7/60=119/120\).
Algebraiczne równanie
Riccatiego sprowadza się do
\[
 p=\frac78+\frac{p}{4(1+p)},
 \qquad
 8p^2-p-7=0.
\]
Ma dwa pierwiastki:
\(p=1\) i \(p=-7/8\).
Pierwszy daje
\(K=(1+p)^{-1}pa=1/4\)
i \(A_{\mathrm{cl}}=1/4\),
więc jest stabilizujący.
Drugi jest ujemny,
a dla niego
\(R+B^*PB=1/8\),
\(K=(-7/8)(1/2)/(1/8)=-7/2\)
oraz \(A_{\mathrm{cl}}=4\).
Nie jest więc stabilizujący.
Zbieżność rekurencji i
jednoznaczność stabilizującego
pierwiastka są zgodne
z Twierdzeniem~\ref{thm:lin24-dare-existence}.
\end{przyklad}
